\begin{abstract}
复杂刚体接触建模是多体动力学中的基础问题。现有互补类硬接触模型能够较好地表达理想刚体非穿透约束，但在任意复杂几何、面接触、分布式接触力矩以及多接触切换场景下，受限于点接触和刚体理想化假设，难以充分刻画接触区域演化及其空间分布作用。现有罚函数或顺应接触模型虽具有更好的工程可实现性与数值连续性，但多数仍停留在点级或少量离散接触点层面，难以准确表达复杂面接触中的接触斑、压力分布及其导致的净接触力与力矩。

针对上述问题，本文提出一种基于符号距离场（Signed Distance Field, SDF）的复杂刚体接触斑—压力场建模方法。该方法以 SDF 作为统一几何表示，首先基于双边距离场定义局部压入与候选接触斑，然后在接触斑上构造局部压力密度，并通过面分布积分计算净接触力与力矩，从而将传统点级顺应接触扩展为面向复杂几何的分布式接触模型。与传统点接触模型相比，本文方法有望更合理地描述复杂刚体之间的面接触、接触中心演化与分布式接触力矩；与完整连续体接触方法相比，该方法不依赖显式体网格和高代价局部变形求解，更适合在多体动力学计算中进行轻量化实现。

本文进一步给出问题定义、接触斑构造方法、压力场闭合关系、数值实现流程以及基准验证框架，为复杂刚体分布式接触建模提供一种面向多体动力学的统一研究路线。
\end{abstract}

\noindent\textbf{关键词：} 多体动力学；符号距离场；接触斑；压力场；顺应接触；复杂刚体接触
